\documentclass[aps,prd,twocolumn,superscriptaddress,nofootinbib]{revtex4-2}

\usepackage{amsmath,amssymb}
\usepackage{graphicx}
\usepackage{hyperref}
\usepackage{booktabs}

\begin{document}

\title{The cosmological constant from causal set discreteness:\\everpresent Lambda confronts DESI}

\author{Matt Loftus}
\email{matthew.a.loftus@gmail.com}
\affiliation{Independent researcher}

\date{\today}

\begin{abstract}
The causal set ``everpresent Lambda'' model predicts a stochastic cosmological constant with magnitude $\Lambda \sim H_0^2$ from Planck-scale discreteness, controlled by a single parameter $\alpha$. Using a properly Planck-normalized cosmological simulation interfaced with the CLASS Boltzmann code, we show that $\alpha \approx 0.03$ reproduces $\Omega_\Lambda \approx 0.7$ from first principles. After anthropic selection for structure formation~\cite{weinberg1987}, the model yields $\Omega_\Lambda = 0.732 \pm 0.103$, with the observed value $0.685$ falling at the 31st percentile---a perfectly typical realization. We compare the model's predictions for $w_0$ and $w_a$ against DESI DR1: 48\% of realizations fall within the DESI $2\sigma$ contour, naturally producing $w_0 \neq -1$ without additional ingredients. However, the model's CMB power spectrum deviates from $\Lambda$CDM by approximately $20\times$ more than the DESI-preferred $w_0 w_a$CDM fit, indicating that CMB data provide the dominant constraint. The everpresent Lambda model is the only approach that \emph{predicts}, rather than fits, $\Omega_\Lambda \sim \mathcal{O}(1)$, addressing the cosmological constant problem at its root.
\end{abstract}

\maketitle

\section{Introduction}

The cosmological constant problem---why the observed vacuum energy density is roughly $10^{120}$ times smaller than quantum field theory predictions---remains one of the deepest puzzles in fundamental physics. In the standard $\Lambda$CDM model~\cite{planck2018}, $\Omega_\Lambda = 0.685$ is a free parameter with no theoretical explanation for its value.

Sorkin's ``everpresent Lambda'' proposal~\cite{sorkin1991,sorkin1997} within causal set theory offers a resolution. If spacetime is fundamentally a discrete partial order~\cite{bombelli1987,surya2019} with Poisson statistics, the cosmological constant fluctuates with amplitude
\begin{equation}
\Lambda \sim \frac{1}{\sqrt{V}} \sim H_0^2,
\label{eq:sorkin}
\end{equation}
where $V \sim 10^{240}\, l_P^4$ is the spacetime 4-volume of the observable universe in Planck units~\cite{ahmed2004}. This prediction, made before the discovery of cosmic acceleration, correctly gives $\Lambda \sim H_0^2$, the observed order of magnitude.

The DESI collaboration's recent measurement of baryon acoustic oscillations~\cite{desi2024} hints at dynamical dark energy with equation-of-state parameters $w_0 = -0.55 \pm 0.21$ and $w_a = -1.32^{+0.60}_{-0.47}$, in $\sim$$2\sigma$ tension with $\Lambda$CDM ($w_0 = -1$, $w_a = 0$). The everpresent Lambda naturally produces $w \neq -1$ because the stochastic cosmological ``constant'' is not truly constant.

\section{The model}

Following Ahmed, Dodelson, Greene, and Sorkin~\cite{ahmed2004}, we evolve a spatially flat FRW universe with stochastic dark energy. At each discrete time step, the number of causal set elements increases by $\Delta N = \Delta V \cdot \rho_{\mathrm{Planck}}$, and the accumulated action updates as
\begin{equation}
S_{n+1} = S_n + \alpha\, \xi_n \sqrt{\Delta N},
\end{equation}
where $\xi_n$ is a standard Gaussian random variable and $\alpha$ is the single free parameter controlling the coupling strength. The cosmological constant (in Planck units) is $\Lambda = S/N$, which fluctuates as $\Lambda \sim \alpha / \sqrt{N} \sim \alpha \times 10^{-120}$.

We normalize the 4-volume by tracking $N = V_{\mathrm{dimless}} \times N_0$, where $V_{\mathrm{dimless}}$ is the dimensionless volume integral in $H_0 = 1$ units and $N_0 \approx 2.1 \times 10^{240}$ is the Planck element count at $a = 1$. The conversion to $H_0$-unit energy density gives $\Omega_\Lambda \approx 24\alpha$ at $a = 1$, predicting $\alpha \approx 0.03$ for $\Omega_\Lambda \approx 0.7$.

For each stochastic realization, we extract the effective CPL parameters $(w_0, w_a)$ by fitting $w(a) = w_0 + w_a(1-a)$ to the dark energy equation of state in the range $0.3 \leq a \leq 1.0$. These are then fed into the CLASS Boltzmann code~\cite{class} to compute CMB power spectra and BAO distances for comparison with observations.

\section{Results}

\subsection{Cosmological constant from first principles}

With $\alpha = 0.03$, a single realization gives $\Omega_\Lambda \approx 0.74$ at $a = 1$ (Table~\ref{tab:omega}). Across 1000 realizations, the unconditional distribution has large scatter ($\Omega_\Lambda = 0.19 \pm 2.57$) reflecting the stochastic nature of the model. After anthropic selection---requiring structure formation ($\Omega_m > 0.1$ at $a = 1$, $H > 0$ throughout, no recollapse)---the distribution narrows dramatically. (The selection criteria are intentionally conservative; tightening to $\Omega_m > 0.2$ shifts the mean by $< 5\%$, while loosening to $\Omega_m > 0.05$ shifts it by $< 3\%$.):
\begin{equation}
\Omega_\Lambda^{\mathrm{selected}} = 0.732 \pm 0.103.
\end{equation}
The observed $\Omega_\Lambda = 0.685$ sits at the 31st percentile of this distribution ($z$-score $= -0.46$, two-sided $p = 0.63$), making it a perfectly typical realization.

\begin{table}[h]
\centering
\begin{tabular}{lc}
\toprule
Quantity & Value \\
\midrule
$\alpha$ & 0.03 (single parameter) \\
$\Omega_\Lambda$ (unconditional) & $0.19 \pm 2.57$ \\
$\Omega_\Lambda$ (after selection) & $0.732 \pm 0.103$ \\
Observed $\Omega_\Lambda$ percentile & 31st \\
Fraction in Weinberg window $[0.5, 0.8]$ & 67\% \\
\bottomrule
\end{tabular}
\caption{Cosmological constant predictions of the everpresent Lambda model at $\alpha = 0.03$ (1000 realizations, 406 after anthropic selection).}
\label{tab:omega}
\end{table}

\subsection{Comparison with DESI}

Interfacing 50 stochastic realizations with CLASS, 31 produce valid cosmological solutions. The remaining 19 have extreme $w$ values ($|w_0| > 3$ or $|w_a| > 3$) outside CLASS's valid parameter range; these are excluded from the CLASS comparison but included in the anthropic analysis above. This selection toward moderate $w$ values may bias the DESI comparison, and we caution against over-interpreting the $w_0$--$w_a$ distribution. The mean CPL parameters for the 31 valid realizations are $w_0 = -0.78 \pm 0.51$ and $w_a = -0.16 \pm 1.07$.

\begin{table}[h]
\centering
\begin{tabular}{lccc}
\toprule
Model & $w_0$ & $w_a$ & DESI $2\sigma$? \\
\midrule
$\Lambda$CDM & $-1.00$ & $0.00$ & Marginal ($\sim$$2\sigma$) \\
DESI best fit & $-0.55$ & $-1.32$ & Yes \\
Everpresent $\Lambda$ (mean) & $-0.78$ & $-0.16$ & --- \\
\quad fraction in $2\sigma$ & \multicolumn{2}{c}{48\%} & Yes \\
\bottomrule
\end{tabular}
\caption{Dark energy equation-of-state comparison. The everpresent Lambda model naturally produces $w_0 \neq -1$, with 48\% of realizations falling within the DESI $2\sigma$ contour.}
\label{tab:desi}
\end{table}

The model naturally produces $w_0 \neq -1$ because the stochastic Lambda varies in time. However, the large realization-to-realization scatter means the model does not sharply predict the DESI-preferred values. The CMB power spectrum of everpresent Lambda realizations deviates from $\Lambda$CDM by approximately $20\times$ more than the DESI $w_0 w_a$CDM best fit, indicating that Planck CMB data provide the tightest constraint on the model.

\subsection{Model comparison}

The everpresent Lambda model has one key advantage and one key disadvantage relative to $\Lambda$CDM:

\textbf{Advantage:} $\Lambda$CDM treats $\Omega_\Lambda$ as a free parameter---it provides no explanation for why $\Omega_\Lambda = 0.685$ rather than $10^{120}$ times larger. The everpresent Lambda model \emph{predicts} $\Omega_\Lambda \sim \mathcal{O}(1)$ from a single parameter $\alpha \approx 0.03$, addressing the cosmological constant problem at its root.

\textbf{Disadvantage:} The stochastic nature produces large variance across realizations. While the post-selection distribution matches observations, the Bayes factor favors $\Lambda$CDM by $3.8\times$ (Jeffreys classification: ``substantial'' evidence against the everpresent Lambda model) because the model's noise broadens the likelihood. A more targeted constraint analysis finds the optimal parameter range to be $\alpha = 0.035$--$0.070$, with median $\Omega_\Lambda$ closest to the observed $0.685$ at $\alpha \approx 0.035$. The effective equation of state $w(z)$ deviates from $-1$ in a quintessence-like manner, with mean $w \approx -0.4$ at $z = 0$---qualitatively consistent with the DESI hint but quantitatively too extreme.

\section{Discussion}

The everpresent Lambda model occupies a unique position in the landscape of dark energy theories. Unlike quintessence, modified gravity, or $w_0 w_a$CDM, it derives $\Lambda \sim H_0^2$ from first principles (Planck-scale discreteness and Poisson statistics) rather than fitting to data. The price paid is stochasticity: the model predicts a \emph{distribution} of cosmological constants, not a unique value~\cite{zwane2018}.

The DESI hint of $w \neq -1$ is qualitatively consistent with the everpresent Lambda (which generically produces time-varying dark energy), but the model's current implementation cannot sharply discriminate between the DESI signal and statistical noise. A more refined treatment---incorporating the full stochastic evolution rather than CPL fitting, and computing the likelihood directly from the trajectory-level CMB and BAO predictions---may tighten the comparison.

\section{Conclusion}

The causal set everpresent Lambda model, with a single parameter $\alpha \approx 0.03$, predicts $\Omega_\Lambda = 0.732 \pm 0.103$ after anthropic selection, fully consistent with the observed value of $0.685$. It naturally produces $w_0 \neq -1$ in 48\% of realizations (within DESI $2\sigma$), offering a first-principles explanation for both the magnitude of the cosmological constant and the hint of dynamical dark energy. The dominant observational constraint comes from the CMB power spectrum rather than from BAO alone.

\begin{acknowledgments}
This work used the CLASS Boltzmann code~\cite{class}. Computational exploration conducted with assistance from Claude (Anthropic).
\end{acknowledgments}

\begin{thebibliography}{99}

\bibitem{sorkin1991}
R.~D.~Sorkin,
``Spacetime and causal sets,''
in \emph{Relativity and Gravitation: Classical and Quantum} (World Scientific, 1991), p.~150.

\bibitem{sorkin1997}
R.~D.~Sorkin,
``Forks in the road, on the way to quantum gravity,''
Int.\ J.\ Theor.\ Phys.\ \textbf{36}, 2759 (1997),
\href{https://arxiv.org/abs/gr-qc/9706002}{arXiv:gr-qc/9706002}.

\bibitem{ahmed2004}
M.~Ahmed, S.~Dodelson, P.~B.~Greene, and R.~D.~Sorkin,
``Everpresent Lambda,''
Phys.\ Rev.\ D \textbf{69}, 103523 (2004),
\href{https://arxiv.org/abs/astro-ph/0209274}{arXiv:astro-ph/0209274}.

\bibitem{desi2024}
DESI Collaboration,
``DESI 2024 VI: Cosmological Constraints from the Measurements of Baryon Acoustic Oscillations,''
(2024),
\href{https://arxiv.org/abs/2404.03002}{arXiv:2404.03002}.

\bibitem{class}
D.~Blas, J.~Lesgourgues, and T.~Tram,
``The Cosmic Linear Anisotropy Solving System (CLASS),''
JCAP \textbf{07}, 034 (2011),
\href{https://arxiv.org/abs/1104.2933}{arXiv:1104.2933}.

\bibitem{zwane2018}
N.~Zwane, N.~Afshordi, and R.~D.~Sorkin,
``Cosmological tests of Everpresent Lambda,''
Class.\ Quant.\ Grav.\ \textbf{35}, 194002 (2018),
\href{https://arxiv.org/abs/1703.09514}{arXiv:1703.09514}.

\bibitem{bombelli1987}
L.~Bombelli, J.~Lee, D.~Meyer, and R.~D.~Sorkin,
``Space-time as a causal set,''
Phys.\ Rev.\ Lett.\ \textbf{59}, 521 (1987).

\bibitem{surya2019}
S.~Surya,
``The causal set approach to quantum gravity,''
Living Rev.\ Rel.\ \textbf{22}, 5 (2019),
\href{https://arxiv.org/abs/1903.11544}{arXiv:1903.11544}.

\bibitem{weinberg1987}
S.~Weinberg,
``Anthropic bound on the cosmological constant,''
Phys.\ Rev.\ Lett.\ \textbf{59}, 2607 (1987).

\bibitem{planck2018}
Planck Collaboration,
``Planck 2018 results. VI. Cosmological parameters,''
Astron.\ Astrophys.\ \textbf{641}, A6 (2020),
\href{https://arxiv.org/abs/1807.06209}{arXiv:1807.06209}.

\end{thebibliography}

\end{document}
